\begin{titlepage}
\centering
\vspace*{2.2cm}

{\LARGE\bfseries Mytheus\par}
\vspace{0.5em}
{\large Personalized Multi-Agent Meeting Orchestration\par}

\vspace{1.6cm}
{\color{rulegrey}\hrule}
\vspace{1.1cm}

{\Large\bfseries A Literature Foundation\par}
\vspace{0.5em}
{\large Twenty five verified papers, what each contributes,\\ and where the field leaves room\par}

\vspace{1.1cm}
{\color{rulegrey}\hrule}
\vspace{2.0cm}

\begin{minipage}{0.82\textwidth}
\small
This document exists to give the Mytheus project a defensible research
footing. Each of the twenty five entries is a real, independently verified
publication. For every one there is a one page treatment covering what the
paper does, how its method actually works, what it reports, the specific
Mytheus component it informs, and the limitations it leaves open.

\vspace{0.7em}
A companion document, \textbf{Research Gaps and Novel Contributions},
aggregates the limitations collected here into the case for what Mytheus can
claim as new.
\end{minipage}

\vfill
{\footnotesize\color{rulegrey}
Companion to the Mytheus architecture documentation set.\par
Every citation in this document was verified against a live API before it was
written up. The protocol is described overleaf.\par}
\vspace{1cm}
\end{titlepage}

\section*{How to read this document}
\addcontentsline{toc}{section}{How to read this document}

The twenty five entries are grouped into eight themes, ordered so that the
document moves from how agents are organised, through what they retrieve, to
how their output is gated before a human sees it. That ordering mirrors the
path a single utterance takes through Mytheus.

Each entry is self contained and roughly one page. If you are looking for
support for a specific design decision rather than reading straight through,
Table \ref{tab:map} routes you from a Mytheus component to the papers that
bear on it.

Two conventions are worth stating. First, the grey block in each entry is the
part written for this project rather than about the paper: it says what to
build, measure or cite, and it is deliberately critical where a published
method would not survive the constraints Mytheus operates under. Second, the
limitations at the foot of each entry are not filler. They are the raw
material for the companion gap analysis, and they were collected with the
question ``what does this paper never test'' rather than ``what did the
authors admit''.

\section*{The verification protocol}

\addcontentsline{toc}{section}{The verification protocol}

Fabricated citations are the characteristic failure of machine assisted
literature work, and a single invented reference is enough to discredit a
dissertation. Every paper in this document was therefore confirmed against a
live authoritative API before it was written up. No citation here rests on
recall.

The protocol ran in three passes:

\begin{enumerate}
\item \textbf{Identifier resolution.} Every candidate carrying an arXiv
identifier was resolved through the arXiv Atom API by identifier list. The
title returned by the API was compared against the claimed title using a
normalised sequence similarity. A candidate was accepted only above a 0.85
match. This pass simultaneously recovered the authoritative title, author
list, publication date and abstract, and those authoritative values, not the
claimed ones, are what appear in this document.

\item \textbf{Bibliographic resolution.} Candidates without an arXiv identifier
were resolved through OpenAlex by title search and through Crossref by DOI,
which also supplied venue and citation counts. Papers published at ACM and
AAAI venues, including the two foundational human computer interaction
references, were confirmed this way.

\item \textbf{Direct retrieval.} Engineering research published outside the
academic system was confirmed by fetching the page and locating the title in
the returned body.
\end{enumerate}

The protocol was validated before use by running it against deliberately
corrupted input. It correctly rejected an invented arXiv identifier, and
correctly rejected a real identifier paired with a wrong title, reporting the
genuine title instead. Of one hundred and thirty candidates gathered, one
hundred and twenty nine resolved. The verification evidence for each paper,
naming the API and the match score, is printed in the header of its entry so
that any reader can repeat the check.

One caveat is stated plainly. Verification establishes that a paper exists,
that its title and authors are correct, and that the link resolves. It does
not certify that every claim summarised in an entry is a faithful reading of
the full paper. Entries were written against the authoritative abstract and
the retrieved paper, and no experimental figure was recorded unless it was
read from the source, but a reader building directly on a specific numerical
result should confirm it in the original.

\newpage
\section*{The map: from research theme to Mytheus component}
\addcontentsline{toc}{section}{The map}

\begin{center}
\begin{tikzpicture}[
  font=\footnotesize,
  stage/.style={draw,rounded corners=2pt,fill=boxgrey,minimum height=0.85cm,
                text width=2.05cm,align=center,inner sep=3pt},
  lbl/.style={text width=2.25cm,align=center,font=\scriptsize\itshape,
              text=accent},
  >={Stealth[round]}
]
\node[stage] (cap) {Capture and window};
\node[stage,right=0.45cm of cap] (agent) {Six agent tier};
\node[stage,right=0.45cm of agent] (peer) {Peer context fill};
\node[stage,right=0.45cm of peer] (ret) {Hybrid retrieval};
\node[stage,right=0.45cm of ret] (gate) {Confidence gate};
\node[stage,right=0.45cm of gate] (conf) {Human confirm};

\draw[->] (cap)--(agent); \draw[->] (agent)--(peer);
\draw[->] (peer)--(ret);  \draw[->] (ret)--(gate);
\draw[->] (gate)--(conf);

\node[lbl,below=0.3cm of cap]   {Streaming and meeting NLP};
\node[lbl,below=0.3cm of agent] {Agent architecture};
\node[lbl,below=0.3cm of peer]  {Multi-agent collaboration};
\node[lbl,below=0.3cm of ret]   {Retrieval and graph memory};
\node[lbl,below=0.3cm of gate]  {Calibration and abstention};
\node[lbl,below=0.3cm of conf]  {Human-AI interaction};

\node[lbl,above=0.3cm of agent, text width=3.2cm] {Personalization};
\node[lbl,above=0.3cm of ret,   text width=3.2cm] {Attribution};
\node[lbl,above=0.3cm of cap,   text width=2.4cm] {Cost control};
\end{tikzpicture}
\end{center}

The vertical labels name the research theme that governs each stage. A paper
placed under a stage is not merely topically adjacent to it: it either
supplies the mechanism Mytheus uses there, or it is the baseline that stage
must be measured against.

{\footnotesize
\begin{longtable}{@{}p{0.5cm}p{6.6cm}p{2.6cm}p{5.4cm}@{}}
\caption{Routing table from Mytheus component to supporting paper.}
\label{tab:map}\\
\toprule
\textbf{\#} & \textbf{Paper} & \textbf{Theme} & \textbf{Mytheus component it bears on} \\
\midrule
\endfirsthead
\toprule
\textbf{\#} & \textbf{Paper} & \textbf{Theme} & \textbf{Mytheus component it bears on} \\
\midrule
\endhead
\bottomrule
\endfoot
1 & AutoGen: Enabling Next-Gen LLM Applications via Multi-Agent Conversation & Agent architecture & BaseAgent lifecycle (spawn/process/terminate) and the NO SILENT COMMITS confirmation contract \\
2 & CAMEL: Communicative Agents for "Mind" Exploration of Large Language Model Society & Agent architecture & ContextGapAgent's agent-to-agent dyadic exchange \\
3 & Why Do Multi-Agent LLM Systems Fail? & Agent architecture & Golden-set evaluation harness, confidence thresholds, and the hallucination-control argument \\
4 & AI Agents That Matter & Agent evaluation & Five-lever cost cascade (relevance pre-gating, batching, LiteLLM tiering, caching, tenant budgets) and the golden-set CI harness \\
5 & Building Effective AI Agents & Agent architecture & BaseAgent lifecycle and the six-agent LangGraph orchestration; relevance pre-gating lever \\
6 & How we built our multi-agent research system & Agent architecture & RAG cite stage / citation contract, and the cost-controlled LLM pipeline \\
7 & Reciprocal rank fusion outperforms condorcet and individual rank learning methods & Retrieval & RAG fuse-and-rerank stage (Reciprocal Rank Fusion, k=60, to top-50) \\
8 & Self-RAG: Learning to Retrieve, Generate, and Critique through Self-Reflection & Retrieval & Uniform confidence thresholds (0.7 / 0.4) across all six agents \\
9 & Proactive Agent: Shifting LLM Agents from Reactive Responses to Active Assistance & Proactive assistance & ProactiveAgent conflict/workload/missing-stakeholder detectors; conflict precision in the golden-set eval harness \\
10 & LaMP: When Large Language Models Meet Personalization & Personalization & PersonalListenerAgent relevance scoring and the role-personalized RAG expand stage via get\_user\_context \\
11 & Personalize Before Retrieve: LLM-based Personalized Query Expansion for User-Centric Retrieval & Personalization & RAG expand stage personalization via get\_user\_context, evaluated against the rerank stage top-10 \\
12 & From Local to Global: A Graph RAG Approach to Query-Focused Summarization & Graph retrieval & Apache AGE organizational-memory property graph + RAG retrieve stage (graph arm) \\
13 & HippoRAG: Neurobiologically Inspired Long-Term Memory for Large Language Models & Graph retrieval & ContextGapAgent RAG fallback tier + RAG fuse/rerank stage \\
14 & Zep: A Temporal Knowledge Graph Architecture for Agent Memory & Graph retrieval & ProactiveAgent conflict detector + Apache AGE organizational-memory graph (FOLLOWS\_FROM lineage) \\
15 & QMSum: A New Benchmark for Query-based Multi-domain Meeting Summarization & Meeting understanding & RAG pipeline (retrieve/rerank/assemble stages) and PersonalListenerAgent relevance scoring \\
16 & Tell me what I need to know: Exploring LLM-based (Personalized) Abstractive Multi-Source Meeting Summarization & Meeting understanding & ContextGapAgent (gap detection plus fallback chain) and PersonalListenerAgent personalization \\
17 & Policies and Evaluation for Online Meeting Summarization & Meeting understanding & 30-second window scheduling, ProactiveAgent live intervention, and the p50<2s / p95<4s agent latency budget \\
18 & On Calibration of Modern Neural Networks & Confidence and abstention & Uniform confidence thresholds (0.7/0.4) and the golden-set eval harness Brier-score calibration check \\
19 & Language Models (Mostly) Know What They Know & Confidence and abstention & ContextGapAgent gap-detection trigger and the below-0.4 discard rule \\
20 & Detecting hallucinations in large language models using semantic entropy & Confidence and abstention & ContextGapAgent silent-on-empty guarantee and the RAG cite stage drop-if-uncitable rule \\
21 & Enabling Large Language Models to Generate Text with Citations & Attribution & Golden-set eval harness (citation validity metric) and the RAG `cite` stage \\
22 & FrugalGPT: How to Use Large Language Models While Reducing Cost and Improving Performance & Cost control & Cost-controlled LLM pipeline (five levers) and tiered model routing via LiteLLM \\
23 & Principles of Mixed-Initiative User Interfaces & Human-AI interaction & ProactiveAgent; uniform confidence thresholds (0.7 surface / 0.4 log / below discard) \\
24 & Guidelines for Human-AI Interaction & Human-AI interaction & NO SILENT COMMITS confirmation contract; ProactiveAgent and QuestionAgent surfacing UX \\
25 & Generative Agents: Interactive Simulacra of Human Behavior & Agent memory & PersonalListenerAgent relevance scoring; confidence thresholds (0.7/0.4) and Brier calibration; ProactiveAgent \\

\end{longtable}
}

\newpage
